\subsection{Estilos Cognitivos}
\label{sec:estilosCognitivos}

Se define como la forma en que el ser humano percibe, recuerda y piensa, o también a la manera en que el ser humano descubre, proceso, almacena, interpreta, transforma y utiliza la información. Estos se adquieren y se modifican gracias a la experiencia \cite{Velez2013}. En la actualidad, hay diversas clasificaciones de estilos cognitivos:  
En \cite{Garcia1989}, se nombran las siguientes:
\begin{itemize}
    \item \textbf{Reflexividad – impulsividad (Kagan):}\ \\
Se refiere a que existen algunas personas que son cuidadosas a la hora de examinar la validez diferencial de varias respuestas alternativas y por lo tanto cometen menos errores; otras son menos cuidadosos y toman decisiones rápidamente sin cuestionar otras alternativas y por ende cometen muchos errores. Los primeros han sido llamados reflexivos y los segundos impulsivos. 

\item \textbf{Dependencia - Independencia del campo (Witkin):}\ \\
Las personas dependientes del campo se caracterizan por ser pasivos, necesitan ser dirigidos y perciben las cosas del contexto de manera estructurada como un todo, siendo difícil el análisis de sus componentes y encontrar relaciones entre sus elementos. Las personas independientes del campo se caracterizan por ser activos no necesitan ser guiados dentro del contexto y son capaces de percibir y analizar las componentes de un todo con mucha facilidad. 

\item \textbf{Convergencia – Divergencia (Getzels y Jackson):}\ \\
Las personas convergentes al enfrentarse a un problema específico solo se centran en solo lo pertinente para dar una solución a ese problema; por otro lado, las personas divergentes al enfrentarse a un problema específico plantean diferentes alternativas para solucionar la problemática, aunque no todas son pertinentes y tienen dificultades al elegir una sola solución.

\item \textbf{Holismo - Serialismo: }\ \\
Las personas holísticas procesan varios elementos a la vez y los organizan empleando analogías y conexiones en aspectos teóricos y conceptuales para crear una unidad compleja; las personas serialistas analizan todos los elementos de un problema en detalle organizándose de manera secuencial y jerárquica.
\end{itemize}

Por su parte en \cite{Pantoja2012} nombran la siguiente clasificación:
\begin{itemize}
\item \textbf{Concreción - Abstracción:}\ \\
 Para el aprendizaje de algo nuevo las personas acuden a ideas concretas o abstractas

\item \textbf{	Independiente – Sensible:}\ \\
La persona tiende a asignarle una organización y estructura propia a la información disponible para realizar una actividad o resolver un problema con independencia de la forma como ha sido presentado, por otra parte, existe la persona que tienden a resolver la actividad o problema manejando la información disponible sin desprenderla del contexto en que ha sido presentada y sin cambiar su estructura y organización inicial. 

\item \textbf{	Adaptación - Innovación:} \ \\
Es la diferencia que tienen las personas para abordar una situación problemática, las adaptativas prefieren hacer las cosas mejor, mientras que las innovadoras prefieren hacer las cosas de una manera diferente.
\end{itemize}
Cabe resaltar que existen diferentes instrumentos de medición que permiten determinar las distintas categorías que posee una persona en cuanto a estilos cognitivos. Entre estos instrumentos se encuentran figuras enmascaradas, test de Emparejamiento de Figuras Familiares, entre otros.

\subsubsection{Estilos Cognitivos en Sistemas Hipermedia Adaptativos}
Anteriormente, en los sistemas Hipermedia adaptativos para el aprendizaje, la adaptación se basaba solo en objetivos o tareas, intereses o nivel de conocimientos del estudiante. Sin embargo, se volvió de gran importancia involucrar los estilos cognitivos en estos sistemas, debido a que estos afectan significativamente el aprendizaje de los estudiantes, puesto que se refieren a cómo los estudiantes procesan y organizan la información \cite{Mampadi2011}.

En este sentido, esta nueva forma de adaptación es consecuencia de un bajo desempeño de los estudiantes frente a una sola forma de enseñanza y, a la necesidad de estrategias adecuadas acordes a la forma en como procesan y perciben la información. Es así como la mayoría de los sistemas hipermedia adaptativos basados en estilos cognitivos, hoy en día, se adaptan principalmente a las opciones de navegación \cite{Ferraro2006}. 
